\subsection{各枝のサイズ}

\begin{definition}
  \label{def:branchSet}
  \lean{}
  \leanok
  集合$B(u, v) \subseteq V$を以下で定める：
  \begin{equation}
    B(u, v) := \set{w}{\textrm{$w$は$u$を通らずに$v$に到達可能}}.
  \end{equation}
  つまり，頂点$u$を削除した後に$v$と繋がっている頂点たちの集合である．
\end{definition}

\begin{lemma}
  \label{lem:branch_ge2_has_neighbor}
  \lean{}
  \leanok
  \uses{}
  $u, v \in V$で，$\# B(u, v) \ge 2$とする．
  このとき，$z \in V$で，以下をすべて満たすものが存在する：
  \begin{equation}
    z \ne v,\ z \ne u,\ C_{vz} \ne 0,\ z \in B(u, v).
  \end{equation}
\end{lemma}

\begin{proof}
  任意の$w \in B(u, v)$に対し，$w = v$ならば，$B(u, v) \subseteq \{v\}$であるので，
  \begin{equation}
    \# B(u, v) \le \# \{v\} < 2
  \end{equation}
  である．
  今，$\# B(u, v) \ge 2$であるので，この対偶より$w \in B(u, v)$で$w \ne v$を満たすものが存在する．
  よって，$u$を通らない$v$から$w$に至る道がある．
  この道の$v$に隣接する頂点を$z$とする．
  すると，この$z$は主張をすべて満たすことがわかる．
\end{proof}

\begin{theorem}
  \label{thm:smallest_branch_eq_one}
  \lean{}
  \leanok
  \uses{}
  $C$はGCMであるとする．
  $S(C)$は正定値であるとする．
  $u \in V$で，$\deg(u) = 3$とする．
  $v_1, v_2, v_3 \in V$は相異なり，$v_1, v_2, v_3 \in N(u),\ \# B(u, v_1) \le \# B(u, v_2) \le \# B(u, v_3)$を満たすとする．
  このとき，以下が成り立つ：
  \begin{equation}
    \# B(u, v_1) = 1.
  \end{equation}
\end{theorem}

\begin{proof}
  $\# B(u, v_1) \ne 1$と仮定して矛盾を導く．

  補題~\ref{lem:branch_ge2_has_neighbor}を$v = v_1, v_2, v_3$として使いたい．
  そこで，各$v_i$に対し，$\# B(u, v_i) \ge 2$を示す．
  $v_1 \in B(u, v_1)$であるから，$\# B(u, v_1) \ne 0$である．
  さらに，仮定より$\# B(u, v_1) \ne 1$である．
  よって，$\# B(u, v_1) \ge 2$が示せた．
  また，$\# B(u, v_1) \le \# B(u, v_2) \le \# B(u, v_3)$であるから$v_2, v_3$についても成り立つ．

  したがって，$w_i \in V$で，以下をすべて満たすものが存在する：
  \begin{equation}
    w_i \ne v_i,\ w_i \ne u,\ C_{v_i w_i} \ne 0,\ w_i \in B(u, v_i) \quad (i = 1, 2, 3).
  \end{equation}
\end{proof}